\Needspace{15\baselineskip}
It remains only for us to give a general method\SourceLeaf{246}%
 for the transformation denoted above by (I). Whatever functions of $x,x',x''$ etc. are taken for $y,y',y''$ etc., it is clear that, at least generally speaking, elimination will allow the quantities $x',x'',x'''$ etc. to be represented as functions of $x,y,y',y''$ etc., whose differentiation will produce $n-1$ equations:
\[
\begin{aligned}
\mathrm{d}x'&=\xi'\mathrm{d}x+\alpha'\mathrm{d}y+\beta'\mathrm{d}y'+\gamma'\mathrm{d}y''+\text{etc.}\\
\mathrm{d}x''&=\xi''\mathrm{d}x+\alpha''\mathrm{d}y+\beta''\mathrm{d}y'+\gamma''\mathrm{d}y''+\text{etc.}\\
\mathrm{d}x'''&=\xi'''\mathrm{d}x+\alpha'''\mathrm{d}y+\beta'''\mathrm{d}y'+\gamma'''\mathrm{d}y''+\text{etc.}\quad\text{etc.}
\end{aligned}
\]
Here, then, the coefficients $\xi',\xi'',\xi'''$ etc., $\alpha',\alpha'',\alpha'''$ etc. are functions of $x,y,y',y''$ etc. In this respect we shall distinguish their partial derivatives with respect to $x$ by enclosing them in parentheses. Evidently those quantities can also be regarded as functions of $x,x',x'',x'''$ etc.; in this respect we shall write the partial derivatives with respect to $x$ without parentheses, so that $\left(\frac{\mathrm{d}\xi'}{\mathrm{d}x}\right)$ must be carefully distinguished from $\frac{\mathrm{d}\xi'}{\mathrm{d}x}$. The same applies to $\left(\frac{\mathrm{d}p}{\mathrm{d}x}\right)$ and $\frac{\mathrm{d}p}{\mathrm{d}x}$ etc. Now, for $\Omega$ to acquire the prescribed form after substituting those values of $\mathrm{d}x',\mathrm{d}x'',\mathrm{d}x'''$ etc., \emph{first}, $\mathrm{d}x$ must evidently disappear, and hence the following condition [1] must hold:
\[
0=p+p'\xi'+p''\xi''+p'''\xi'''+\text{etc.}
\]
Furthermore, the coefficients of $\mathrm{d}y,\mathrm{d}y',\mathrm{d}y''$ etc., namely
\[
\begin{aligned}
p'\alpha'+p''\alpha''+p'''\alpha'''+\text{etc.}&=A\\
p'\beta'+p''\beta''+p'''\beta'''+\text{etc.}&=B\\
p'\gamma'+p''\gamma''+p'''\gamma'''+\text{etc.}&=C,\quad\text{etc.}
\end{aligned}
\]
are to acquire the form $\lambda q,\lambda q',\lambda q''$ etc., so that $q,q',q''$ etc. become functions only of $y,y',y''$ etc. For this to happen, we must, \emph{second}, have [2]:
\[
\frac{1}{A}\cdot\left(\frac{\mathrm{d}A}{\mathrm{d}x}\right)
=\frac{1}{B}\cdot\left(\frac{\mathrm{d}B}{\mathrm{d}x}\right)
=\frac{1}{C}\cdot\left(\frac{\mathrm{d}C}{\mathrm{d}x}\right)\text{etc.}
=\frac{1}{\lambda}\cdot\left(\frac{\mathrm{d}\lambda}{\mathrm{d}x}\right)
\]
Now,
\[
\begin{aligned}
\left(\frac{\mathrm{d}A}{\mathrm{d}x}\right)={}&
p'\left(\frac{\mathrm{d}\alpha'}{\mathrm{d}x}\right)
+p''\left(\frac{\mathrm{d}\alpha''}{\mathrm{d}x}\right)
+p'''\left(\frac{\mathrm{d}\alpha'''}{\mathrm{d}x}\right)+\text{etc.}\\
&+\alpha'\left(\frac{\mathrm{d}p'}{\mathrm{d}x}\right)
+\alpha''\left(\frac{\mathrm{d}p''}{\mathrm{d}x}\right)
+\alpha'''\left(\frac{\mathrm{d}p'''}{\mathrm{d}x}\right)+\text{etc.}
\end{aligned}
\]
\SourceLeaf{247}
\Needspace{20\baselineskip}
If we substitute here
\[
\left(\frac{\mathrm{d}\alpha'}{\mathrm{d}x}\right)=\frac{\mathrm{d}\xi'}{\mathrm{d}y},\qquad
\left(\frac{\mathrm{d}\alpha''}{\mathrm{d}x}\right)=\frac{\mathrm{d}\xi''}{\mathrm{d}y}.\quad\text{etc.}
\]
and subtract
\[
\begin{aligned}
0={}&\frac{\mathrm{d}p}{\mathrm{d}y}
+\xi'\cdot\frac{\mathrm{d}p'}{\mathrm{d}y}
+\xi''\cdot\frac{\mathrm{d}p''}{\mathrm{d}y}+\text{etc.}\\
&+p'\cdot\frac{\mathrm{d}\xi'}{\mathrm{d}y}
+p''\cdot\frac{\mathrm{d}\xi''}{\mathrm{d}y}+\text{etc.}
\end{aligned}
\]
which is the equation obtained by differentiating [1] with respect to $y$, then
\[
\left(\frac{\mathrm{d}A}{\mathrm{d}x}\right)=
\left\{\begin{aligned}
&\phantom{-\frac{\mathrm{d}p}{\mathrm{d}y}}+\alpha'\cdot\left(\frac{\mathrm{d}p'}{\mathrm{d}x}\right)
+\alpha''\cdot\left(\frac{\mathrm{d}p''}{\mathrm{d}x}\right)+\text{etc.}\\
&-\frac{\mathrm{d}p}{\mathrm{d}y}-\xi'\cdot\frac{\mathrm{d}p'}{\mathrm{d}y}
-\xi''\cdot\frac{\mathrm{d}p''}{\mathrm{d}y}-\text{etc.}
\end{aligned}\right.
\]
Since, moreover,
\[
\begin{aligned}
\left(\frac{\mathrm{d}p'}{\mathrm{d}x}\right)&=\frac{\mathrm{d}p'}{\mathrm{d}x}
+\xi'\cdot\frac{\mathrm{d}p'}{\mathrm{d}x'}
+\xi''\cdot\frac{\mathrm{d}p'}{\mathrm{d}x''}+\text{etc.}\\
\left(\frac{\mathrm{d}p''}{\mathrm{d}x}\right)&=\frac{\mathrm{d}p''}{\mathrm{d}x}
+\xi'\cdot\frac{\mathrm{d}p''}{\mathrm{d}x'}
+\xi''\cdot\frac{\mathrm{d}p''}{\mathrm{d}x''}+\text{etc.},\quad\text{etc.}
\end{aligned}
\]
\[
\begin{aligned}
\frac{\mathrm{d}p}{\mathrm{d}y}&=\alpha'\cdot\frac{\mathrm{d}p}{\mathrm{d}x'}
+\alpha''\cdot\frac{\mathrm{d}p}{\mathrm{d}x''}
+\alpha'''\cdot\frac{\mathrm{d}p}{\mathrm{d}x'''}+\text{etc.}\\
\frac{\mathrm{d}p'}{\mathrm{d}y}&=\alpha'\cdot\frac{\mathrm{d}p'}{\mathrm{d}x'}
+\alpha''\cdot\frac{\mathrm{d}p'}{\mathrm{d}x''}
+\alpha'''\cdot\frac{\mathrm{d}p'}{\mathrm{d}x'''}+\text{etc.},\quad\text{etc.}
\end{aligned}
\]
these substitutions give
\[
\left(\frac{\mathrm{d}A}{\mathrm{d}x}\right)=k'\alpha'+k''\alpha''+k'''\alpha'''+\text{etc.}
\]
where
\[
\begin{array}{r@{}l@{}l@{}l@{}l@{}l}
k'={}&(1,0)&\quad *\quad&+(1,2)\xi''&+(1,3)\xi'''&+\text{etc.}\\
k''={}&(2,0)&+(2,1)\xi'&\quad *\quad&+(2,3)\xi'''&+\text{etc.}\\
k'''={}&(3,0)&+(3,1)\xi'&+(3.2)\xi''-&\quad *\quad&+\text{etc.}
\end{array}
\]
etc., if, for brevity, we generally denote
\[
\frac{\mathrm{d}p^{(\mu)}}{\mathrm{d}x^{(\nu)}}
-\frac{\mathrm{d}p^{(\nu)}}{\mathrm{d}x^{(\mu)}}
\]
by $(\mu,\nu)$, so that in general $(\mu,\mu)=0$ and $(\nu,\mu)=-(\mu.\nu)$.
It is also easy to see that
\[
\begin{aligned}
\left(\frac{\mathrm{d}B}{\mathrm{d}x}\right)&=k'\beta'+k''\beta''+k'''\beta'''+\text{etc.}\\
\left(\frac{\mathrm{d}C}{\mathrm{d}x}\right)&=k'\gamma'+k''\gamma''+k'''\gamma'''+\text{etc.}
\end{aligned}
\]
\SourceLeaf{248}
etc., and consequently equations [2] will be satisfied if \(k',\,k'',\,k'''\) etc. become proportional, respectively, to the quantities \(p'.\,p'',\,p'''\) etc. If we also set
\[
k=\quad *\quad +(0,1)\xi'+(0,2)\xi''+(0,3)\xi'''+\text{etc.}
\]
we have the identity
\[
0=k+k'\xi'+k''\xi''+k'''\xi'''+\text{etc.}
\]
from which, combined with [1], it follows readily that \(k\) must also be proportional to \(p\). This latter proportionality can replace equation [1]. With the aid of the \(n-1\) equations
\[
\frac{k}{p}=\frac{k'}{p'}=\frac{k''}{p''}=\frac{k'''}{p'''}=\text{etc.}
\]
the hitherto unknown functions \(\xi',\,\xi'',\,\xi'''\) etc., whose number is likewise \(n-1\), can now be determined. Closer consideration shows, however, that this determination can be carried out only for even values of \(n\). For odd \(n\), as soon as all but one of the quantities \(\xi'.\,\xi'',\,\xi'''\) etc. have been eliminated, this last one will always disappear of itself, leaving only a \emph{equation expressing a condition} between \(p,\,p',\,p''\) etc. This circumstance is the reason why transformation (I) must be restricted to even values of \(n\).

The conditions for transformation I have thus now been reduced to the requirement that the partial derivatives \(\bigl(\frac{\mathrm{d}x'}{\mathrm{d}x}\bigr),\ \bigl(\frac{\mathrm{d}x''}{\mathrm{d}x}\bigr),\ \bigl(\frac{\mathrm{d}x'''}{\mathrm{d}x}\bigr)\) etc., insofar as \(x',\,x'',\,x'''\) etc. are regarded as functions of \(x,\,y,\,y',\,y''\) etc., take the values now represented as known functions of \(x,\,x',\,x''\) etc., namely \(\xi',\xi'',\xi'''\) etc. This can also be expressed as follows: insofar as \(y,\,y',\,y''\) etc. are regarded as constant, and hence \(x',\,x'',\,x'''\) etc. only as functions of the variable \(x\), the following \(n-1\) differential equations must be satisfied:
\[
\mathrm{d}x'=\xi'\mathrm{d}x,\qquad
\mathrm{d}x''=\xi''\mathrm{d}x,\qquad
\mathrm{d}x'''=\xi'''\mathrm{d}x,\quad\text{etc.}
\]
or, if one prefers the original equations from whose combination these actually arose, the following:
\SourceLeaf{249}
\[
\begin{aligned}
&\quad *\qquad +\frac{(0,1)}{p}\mathbin{.}\mathrm{d}x'
 +\frac{(0,2)}{p}\mathbin{.}\mathrm{d}x''
 +\frac{(0,3)}{p}\mathbin{.}\mathrm{d}x'''+\text{etc.}\\
&=\frac{(1,0)}{p'}\mathbin{.}\mathrm{d}x\qquad *\quad
 +\frac{(1,2)}{p'}\mathbin{.}\mathrm{d}x''
 +\frac{(1,3)}{p'}\,\mathrm{d}x'''+\text{etc.}\\
&=\frac{(2,0)}{p''}\mathbin{.}\mathrm{d}x
 +\frac{(2,1)}{p''}\mathbin{.}\mathrm{d}x'\qquad *\quad
 +\frac{(2,3)}{p''}\mathbin{.}\mathrm{d}x'''+\text{etc.}\\
&=\text{etc.}
\end{aligned}
\]
\Needspace{22\baselineskip}
The integration of these equations belongs, however, to ordinary integral calculus, and is presupposed here. Generally speaking, it will contain \(n-1\) independent constants \(H,\,H',\,H'',\,H'''\) etc., which appear as given functions of \(x',\,x'',\,x'''\) etc., so that one obtains \(n-1\) finite equations
\[
H=X,\qquad H'=X',\qquad H''=X'',\quad\text{etc.}
\]
if \(X,\,X',\,X''\) etc. represent these functions. This analysis thus shows that the prescribed conditions will be satisfied if these very functions are chosen for \(y,\,y',\,y''\) etc., or if
\[
y=X,\qquad y'=X',\qquad y''=X'',\quad\text{etc.}
\]
is set.

We wish to add one further remark here. We have deliberately stated that the integration gives, \emph{generally speaking}, \(n-1\) independent constants. This no longer holds in special cases: namely, if the coefficients \(p,\,p',\,p''\) etc. are such that the above \(n-1\) differential equations are not independent, but one can already be derived by combining the others. In this case, the determination of \(\xi',\,\xi'',\,\xi'''\) etc. by elimination will no longer be possible either. This case ought properly to be treated separately as an exception. Here, however, we are all the more content with a \emph{brief indication} of how to proceed, since the author has not touched upon it. To apply what precedes, one need only replace the one differential equation above that is already contained in the others by any other arbitrary linear equation between \(\mathrm{d}x,\,\mathrm{d}x',\,\mathrm{d}x''\) etc. not already contained in them. It will always be most convenient to set one of these differentials \(=0\), or to treat one of the variables \(x,\,x',\,x''\) etc.\SourceLeaf{250}
 as constant, for example \(x\), unless \(\mathrm{d}x=0\) already happens to follow from the existing equations. In this case, one of the integral equations will be
\[
H=x
\]
and, if one sets \(y=x\), it can be shown that in the transformed expression for \(\Omega\) the relation \(q=0\) always holds automatically.

\begin{center}\rule{0.25\linewidth}{0.3pt}\end{center}


